 \vspace{4cm}
 At this point, we have all the necessary tools to describe the results
 of this thesis.
 The first one, reported in Chapter 5, is a gravitational wave-like propa
gation on the top of an acoustic metric. Notably, in the past only static
 solutions of GR were designed in analogue models. Here,
 we present a system where analogue gravitational waves propagate sat
isfying Einstein’s equations.
 The second result, in Chapter 6, consists in designing a system in which
 an acoustic horizon is excited by a gravitational wave-like perturbation.
 To this purpose, we choose a cylindrical acoustic black hole, where we
 introduce an appropriate extension in the geometry of the previously
 found gravitational wave-like perturbation.


\newpage


\chapter{Simulation setup and implementation}
\label{Simulation_setup}

\section{Lensed BBH population with LER}
% Explicar qué es LER y para qué lo usas
% Parámetros: npool, event_type BBH, ifos L1 H1
% PSD: T2500310-v2_O5a_strain_psd (O5 design sensitivity)
% Waveform: IMRPhenomXPHM, spin precession
% Criterio de selección: 500 eventos con SNR >= 10 en ambas imágenes

\section{Injected parameter distributions}
% Plots de m1, m2 con crosses de los valores inyectados
% Distribución de masas, spins, distancias
% Histogram SNR de L1, H1 y joint para ambas imágenes
% Plots LER de los parámetros de lensing: time delay, magnification ratio


\chapter{Parameter estimation pipeline}
\label{Parameter_estimation_pipeline}

\section{Lensed waveform generation with Hanabi}
% Qué es Hanabi y cómo genera señales lensadas
% strongly_lensed_BBH_waveform con IMRPhenomXPNR
% image_type (Morse phase), relative_magnification
% Cómo se conecta LER → Hanabi → Bilby

\section{Bilby configuration and joint analysis}
% Estructura del ini file: label, outdir, request-cpus=128
% Sampler: dynesty, nlive=1000, naccept=60, acceptance-walk
% n-triggers=2, trigger-ini-files (image1 + image2)
% common-parameters: chirp_mass, mass_ratio, spins, ra, dec, psi, phase
% lensing-prior-dict: relative_magnification, image_type DiscreteUniform
% Por qué se hace análisis conjunto (joint) y no individual

\section{Automation and infrastructure}
% Cómo automatizas la creación de ini files para 500 inyecciones
% Submission a Picasso con SLURM
% Gestión de outputs y JSON results

\section{Residuals and parameter recovery}
% Violin plots  
% residuals (parámetros recuperados vs inyectados)
% Scatter plots
\subsection{SNR comparison}
% Comparación entre el SNR calculado por LER y el recuperado por Bilby
% Validación del setup}

\chapter{Sky localization}
\label{Skyloc}

\section{Sky map generation}
% ligo-skymap-from-samples: genera el skymap.fits a partir 
% de los posterior samples de Bilby
% --disable-distance-map, --trials 1
% Se hace para cada imagen individual y para el joint

\section{Sky area and constellation analysis}
% ligo-skymap-stats: extrae el área al 90% de credibilidad
% ligo-skymap-constellations: asigna la región a constelaciones
% Plot SNR vs sky area 90%
% Plot SNR vs constellation

\section{Comparison between single image and joint analysis}
% PDF y CDF del sky area comparando:
%   - imagen individual (la mejor, mayor SNR)
%   - análisis joint de las dos imágenes
% Mejora cuantitativa en la localización



FUTURE WOOOORKKK


To summarize, in this chapter we have developed the description of an acoustic
 cylindrical black hole perturbed by a gravitational wave-like perturbation. This
 perturbation is introduced as the cylindrical extension of the analogue gravitational
 wave found in Chapter 5. After having illustrated the designed physical system,
 we have studied how the acoustic horizon is perturbed by the impinging analogue
 gravitational wave, computing also its generators. This analogue model is important
 since it allows to study how an acoustic horizon responds to a gravitational wave-like
 perturbation.